\section{Introduction}
  \label{sec:intro}

  This paper develops a local thermodynamic interpretation of spectral geometry.
  The central object is the projective entropy functional
  \begin{equation}
    S_\Pi[g] \equiv \frac12 \log \det{}' A_g ,
  \end{equation}
  for a Laplace-type operator \(A_g\) on a smooth manifold \((M,g)\).

  We show that the local heat-kernel structure naturally defines a local spectral energy density \(u(x;t)\).
  Its curvature expansion induces a local inverse temperature field \(\beta(x)\) whose leading correction is geometric:
  \begin{equation}
    \beta(x)^{-1} = \beta_0^{-1} + \frac16 R(x) + \ldots
    \label{eq:beta-curvature}
  \end{equation}
  This provides a direct parallel with local Unruh-type constructions, in which temperature acts as a local Lagrange
  multiplier encoding the cost of maintaining a constrained description.

  We then formulate a local spectral first law,
  \begin{equation}
    \delta S_\Pi = \int_{\Sigma} \beta(x)^{-1}\,\delta E_\Pi(x)\, d^3x ,
    \label{eq:local-first-law}
  \end{equation}
  and use it to derive an Einstein response as a condition of local spectral equilibrium.
  In this sense, the field equation
  \(
  G_{\mu\nu} = 8\pi G\,T_{\mu\nu}
  \)
  appears as a thermodynamic compatibility constraint, rather than a fundamental postulate.
